% Source: https://artofproblemsolving.com/community/c3h1603057_2017_mathcounts_nationals_12infinite_geometric_series
% Stable identity: aops-topic-1603057-q01
The successive changes are $+12$, $-6$, $+3$, $-\dfrac{3}{2}$, $\dots$: a geometric sequence with first term $12$ and ratio $-\dfrac{1}{2}$. After $100$ cycles (each an addition followed by a drain, so $200$ changes) the tank holds
\[
12\cdot\dfrac{1-\left(-\dfrac{1}{2}\right)^{200}}{1-\left(-\dfrac{1}{2}\right)}=8\left(1-\dfrac{1}{2^{200}}\right)
\]
gallons, which differs from $8$ by far less than $\dfrac{1}{2}$ (the infinite sum is $\dfrac{12}{1+\dfrac{1}{2}}=8$). So the tank holds $\boxed{8}$ gallons to the nearest whole number. (Counting a cycle as a single step instead changes nothing: after any number of steps beyond the first few the amount is within $\dfrac{1}{2}$ of $8$.)
